\section{Discussion and open questions}
\label{sec:discussion}

\subsection{What the results establish}
\label{sec:discussion-summary}

The four outputs of Section~\ref{sec:intro} have different complexities.
The results locate boundaries for polynomial problems with explicit rational
data, under the convexity hypotheses stated in each theorem, and give related
arithmetic interfaces for nonconvex quadratic constraints and cone problems
with explicit algebraic data.

Values become points when an effective growth modulus is available. Global
convexity supplies one at every degree, with exponent $1/D$ and a constant of
polynomial bit length, so a fixed optimizer can be approximated in time
$\poly(L,D,q)$ (Theorem~\ref{thm:global-point}). Convexity on the feasible set
supplies one for cubics (Theorem~\ref{thm:cubic-point}) but not for quartics,
where a constant-accuracy point already decides Square Root Sum or $\PosSLP$
(Section~\ref{sec:points}). For cubics on product boxes that are convex in
the residual variables but possibly nonconvex in a core of dimension $k$, one
random tilt of the core gives full point output for the tilted instance,
correct for every draw and every precision, with expected work
$k^{O(k)}(1+\Lambda/\sigma)^k\poly(L+q)$ (Theorem~\ref{thm:recourse-cubic}).

Points become exact decisions when a separation bound can be reached. For
explicit globally strongly convex objectives of unary degree, every exact
predicate at the minimizer reduces to one $\PosSLP$ instance
(Theorem~\ref{thm:exact-upper}). Order comparisons of the minimum value are
$\PosSLP$-hard already for quartics with a checked strict convexity
certificate, even when the minimum is promised to be nonzero, and order
comparisons of a minimizer coordinate remain hard when the minimizer is
rational and the minimum is known to be zero
(Theorems~\ref{thm:quartic-complete} and~\ref{thm:rational-optimizer}). In
particular, deciding whether a certified quartic inequality $f(x)\le0$ has a
real solution is $\PosSLP$-complete. Linear constraints add the search for
the active face. Over arbitrary rational polyhedra the comparison lies in
$\mathrm{UP}^{\PosSLP}\cap\mathrm{coUP}^{\PosSLP}$, and our methods remove the
search under structural assumptions, a small nonlinear dimension, or
randomization; integer variables add a discrete search that needs only
ordinary precision (Section~\ref{sec:constraints}).

Exact descriptions have costs that are not determined by the cost of exact
decisions. The coordinates of points isolated by rational convex polynomial
inequalities are characterized (Theorem~\ref{thm:singleton-field}). Their
degree, the height of rational optimizers and witnesses, the size of interior
Gram matrices, and the degree of the least coefficient field of a
sum-of-squares certificate can all grow exponentially with the dimension for
certified inputs whose length is polynomial in the dimension
(Sections~\ref{sec:algebraic}--\ref{sec:fields-certificates}). For the
families that prove these lower bounds, the paper also gives polynomial-size
descriptions of the same objects in compact formats: repeated-squaring
circuits for the long rational minimizers, shared rational circuits for
strictly feasible points and for interior Gram matrices, a root circuit with
one real root gate for the coordinates of the cyclic family, and root
circuits with fifth-root gates, or auxiliary variables for the coordinates of
the minimizer, for the certificates of the tower family. A rational circuit
has rational values, so it cannot describe an irrational coordinate; root
gates, algebraic representations, or implicit descriptions are needed there.
Lower bounds on output size therefore depend on the format and should be
stated for a specific one.

Common-field representations also let exact outputs serve as inputs to a
later computation. Appendix~\ref{app:further-arithmetic} bounds the degree
and total length of feasible witnesses, finite infima, and attained
minimum-norm optimizers for rational quadratic constraints whose Hessians
span a space of dimension $h$, without assuming convexity. It separately
gives exact cone feasibility and witness recovery over one explicitly
represented real number field, with bounded integer dimension and few
independent continuous Hessians of the squared cone residuals. These are
specified algebraic input and output contracts; they do not turn the
nonconvex problems into ordinary polynomial-time optimization in unrestricted
dimension.

\subsection{Consequences for exact and verified optimization}
\label{sec:discussion-consequences}

The results bear on the design of exact and proof-producing solvers in the
following ways. None of them is a running-time claim about existing
software.

\paragraph{Bounds and pruning.}
For globally convex relaxations of fixed or unary degree, certified value
enclosures cost polynomial time, so a pruning test with a positive tolerance
can be carried out in polynomial time. Exact pruning
decisions, such as whether a relaxation bound is strictly below an incumbent
value, are $\PosSLP$-hard in the worst case even for unconstrained quartics
with a checked strict convexity certificate. The hardness is a statement about
how close to zero a nonzero optimal value can be. If there were a polynomial
$s$ such that every nonzero minimum in this class satisfied
$\abs{f_*}\ge2^{-s(L)}$, then approximating $f_*$ to $s(L)+1$ bits would decide
its sign in polynomial time, and $\PosSLP$ would be in P. Exact methods that
refine numerical approximations must therefore expect to need precision that
is not bounded by any fixed polynomial in $L$ on some inputs, unless $\PosSLP$
is in P. In this class, Theorem~\ref{thm:exact-upper} shows that $2^{e(L)}$
bits of precision always suffice, and its Newton circuits reach that
precision with polynomially many gates.

\paragraph{Points and active sets.}
Branching rules and rounding heuristics use points. For globally convex
relaxations of fixed or unary degree, a fixed approximate optimizer is
available in polynomial time. Relaxations built on a box are often convex
only on that box. For quartics of this kind, no deterministic polynomial-time
algorithm returns a point within distance $1/4$ of the minimizer on every
instance with a unique minimizer, unless Square Root Sum and $\PosSLP$ lie in
P (Theorem~\ref{thm:points-quartic-lower}); the hard instances are
constructed for that theorem and are not claimed to arise from a particular
relaxation scheme. Whether a bound is active at the relaxation optimizer is
an exact decision: it is Square Root Sum-hard for strongly convex cubics on a
box, although their optimizers are easy to approximate. For a globally
strongly convex objective of unary degree, a supplied lower bound on the
nonzero slacks lets ordinary approximation recover the active set
(Proposition~\ref{prop:upper-slack-gap}); without such a bound, an active set
read off from an approximate solution carries no guarantee. These
statements concern the active set of the relaxation optimizer. Using it to
fix variables of the original mixed-integer problem needs an additional
optimization argument, which this paper does not provide.

\paragraph{Integer search and exact selection.}
For mixed-integer problems whose quartic objective is strongly convex in all
variables jointly, candidate lists that contain every optimal integer block
can be built with ordinary arithmetic in time that is fixed-parameter
tractable in the number of integer variables, and exact arithmetic is needed
only to select among the candidates. This
separation suggests performing the discrete search numerically and reserving
exact or algebraic arithmetic for the final comparisons. The binary
extraction result shows that the final comparisons cannot be avoided in
general, even when the optimal value is known.

\paragraph{Certificates for independent checking.}
A rational strict convexity certificate does not imply a rational
sum-of-squares certificate of optimality: the ternary quartic $F_\star$ of
Example~\ref{ex:fields-ternary} has minimum zero and neither a rational
sum-of-squares decomposition nor a rational positive semidefinite Gram
matrix. The alternatives have different costs. A rational certificate for
$F_\star+\eta$ exists for every rational $\eta>0$, but our results give no
bound on its size as $\eta$ tends to zero. The common denominator
$1+\norm{x}^2$ gives a short exact rational certificate for $F_\star$
(Example~\ref{ex:fields-ternary-multiplier}) and, after a polynomial-size
increase of a scale parameter, for the tower family, whose polynomial
certificates need coefficients of degree $5^k$
(Section~\ref{sec:fields-certificates}). Auxiliary variables for the
irrational coordinates, together with rational quadratic equations that they
satisfy, also give polynomial-size rational certificates for the tower
family. When the minimum
is positive, Theorem~\ref{thm:circuit-gram} gives a polynomial-size shared
rational circuit for an interior Gram matrix, but checking that the matrix is
positive definite is as hard as deciding the sign of the minimum. Imposing
additional structure on a certificate, such as separation by blocks of
variables or a fixed power of $1+\norm{x}^2$ as multiplier, can make it fail to
exist or force it to be long, while the unrestricted or adapted certificate
remains short.

\paragraph{Reporting exact optimizers.}
For the constrained strongly convex problems of
Section~\ref{sec:constraints}, an active face and a polynomial system with a
unique real solution on that face provide an implicit exact description.
Expanded rational coordinates and coefficient lists of minimal polynomials
can be exponentially long in the dimension even for well-conditioned instances
with small coefficients.

\paragraph{Passing algebraic data between problems.}
A cone problem over rational data can acquire algebraic coefficients when an
exact value from a preceding computation is used as a threshold.
Appendix~\ref{app:further-arithmetic} handles known algebraic thresholds by
placing all data in one specified real number field. With fixed integer
dimension $t$ and a fixed dimension $h$ for the span of the continuous
Hessians of the squared cone residuals, its feasibility and
witness algorithms run in ordinary polynomial time and require neither a
bounding box nor a Slater condition. The continuous witness is represented
in the joint field generated by the input field and its coordinates. This
charges the cost of selecting the real embedding and writing a common-field
output, which a coordinatewise degree estimate alone would not control.

\subsection{Open questions}
\label{sec:discussion-open}

The following questions are not settled by the results above. No theorem of
this paper depends on an answer to any of them.

\begin{enumerate}[label=\arabic*.,leftmargin=*]
\item \emph{Arbitrary polyhedra.} Is exact comparison for globally strongly
  convex quartics on arbitrary rational polyhedra in $\mathrm P^{\PosSLP}$,
  equivalently reducible to one $\PosSLP$ instance? It lies in
  $\mathrm{UP}^{\PosSLP}\cap\mathrm{coUP}^{\PosSLP}$, and our methods remove
  the search under structure, small nonlinear dimension, or randomization.
  The constrained Newton recurrence converges without knowledge of the active
  set. What our approach lacks is an exact algorithm for convex quadratic
  programs whose number of arithmetic operations is polynomial independently
  of the encoding length of the Hessian, because the Hessians of the Newton
  subproblems are given by circuits. Bit-polynomial quadratic programming
  algorithms, smooth penalties, and barrier methods with ordinary accuracy
  bounds do not supply this interface (Section~\ref{sec:constraints-open}).
\item \emph{Equality.} Is deciding $h(p)=0$ $\PosSLP$-hard in the class of
  Theorem~\ref{thm:quartic-complete}, or is it easier? Our lower bounds
  concern order relations only. Deciding whether an integer circuit computes
  zero, the problem EquSLP, lies in coRP by a result of Sch\"onhage that
  Allender et al.\ record~\cite[Section~2]{AllenderEtAl2009}. This fact
  concerns zero tests of integer circuits, not signs, and it does not apply
  directly here: $h(p)$ is the limit of the Newton circuits rather than the
  value of one, and the reduction of Theorem~\ref{thm:exact-upper} produces a
  sign test.
\item \emph{Degenerate convex quartics.} Is there an upper bound relative to
  $\PosSLP$ for exact comparisons at the minimum-norm optimizer of a globally
  convex quartic that is not strongly convex? If the Hessian vanishes entirely
  at an optimizer, the optimizer set is a rational affine space computable by
  linear algebra. For partially degenerate optimizers, Newton's method and its
  simple accelerations fail to converge quadratically, and the Hessian kernel
  at the optimizer need not contain a nonzero rational vector
  (Appendix~\ref{app:boundaries}). These examples rule out particular proof
  routes, not an upper bound.
\item \emph{Rational sums of squares at minimum zero.} In the class of
  Theorem~\ref{thm:quartic-complete}, membership in the cone of rational sums
  of squares is $\PosSLP$-hard. Is it in $\mathrm P^{\PosSLP}$? At a supplied
  zero of the tower family it reduces to rational linear algebra and one
  semidefiniteness test (Section~\ref{sec:fields}), but in general the zero
  is unknown and the answer depends on its coordinate field.
\item \emph{Size of all Gram matrices.} Is there a family of short certified
  quartics with positive minimum for which every rational positive
  semidefinite polynomial Gram matrix has superpolynomial bit length? Our
  lower bounds concern interior, maximal-rank, exposing, or block-separated
  certificates. The families in this paper have short singular rational Gram
  matrices, and the determinant argument for interior Gram matrices does not
  apply to singular ones.
\item \emph{Degree maxima.} For globally convex rational quartics that are
  rational sums of squares and have a unique zero with positive definite
  Hessian, the largest possible degree of the coordinate field is known for
  $n\le5$. For $n\ge6$ it lies between $(2^{n+1}-(-1)^{n+1})/3$ and $2^n-5$;
  for $n=6$, between $43$ and $59$. Without the sum-of-squares hypothesis
  Section~\ref{sec:algebraic} proves the upper bound $2\cdot3^{n-1}-1$, and the
  exact maxima are not known for $n\ge3$.
\item \emph{Domain-convex cubics.} The optimal error exponent for cubics that
  are convex on a polytope lies between $1/4$, proved here, and $1/3$, forced
  by $t^3$ on $[0,1]$.
\item \emph{Recourse beyond product boxes.} Can full point output under core
  noise be obtained for residual-convex cubics on polytopes that couple core
  and residual variables, without additional structure such as a supplied
  convexifier for the core? For residual quartics, a guarantee of
  the form of Section~\ref{sec:recourse} would give a Las Vegas expected
  polynomial-time algorithm for $\PosSLP$ (Remark~\ref{rem:recourse-quartic}),
  so further structure is needed.
\item \emph{Input models.} Our point theorem and exact upper bound are
  polynomial in the numerical degree. For sparse polynomials with
  exponents written in binary, or for objectives given by arithmetic
  circuits, the value, point, and exact problems need new arguments.
\end{enumerate}

Finally, every reduction to one $\PosSLP$ instance in this paper would become
an ordinary polynomial-time algorithm if $\PosSLP$ were in P, whereas the
results stated with $\mathrm{UP}^{\PosSLP}\cap\mathrm{coUP}^{\PosSLP}$ or with
Las Vegas oracle algorithms would yield the corresponding classes without the
oracle. The size lower bounds for expanded outputs and the field lower bounds
for certificates do not depend on the complexity of $\PosSLP$.
